%! Complex Numbers and the Library of Parent Functions — Unit 2, Lesson 2.1 — AP Practice
% EXTRA CREDIT — optional, exactly TWO pages, placed between the notes and the graded homework.
% Shaped like the AP Precalculus exam: page 1 is Section I, five multiple-choice items with four
% options (A)–(D); page 2 is Section II, one multi-part free-response question on a pre-drawn
% graph. Its context (a machined part's length error) is used nowhere else in the lesson.
% Distractors are the lesson's errors on purpose: MC 1 treats i^2 as +1 or drops the i^2 term,
% MC 2 multiplies under one root (sqrt(-4)sqrt(-25) = sqrt(100)), MC 3 forgets the second root
% or decides "no solutions", MC 4 confuses cube root with cube or square root, MC 5 puts the
% -(x - 5) rule on the wrong side of the break (the lesson's target misconception).
% Answers: C, B, D, A, B.
%
% PAGE LOCKSTEP with the other half of the pair: the key marks the correct option in its
% LABEL (no text added to the line) and prose answers are \writespace{H}{…}, fixed height both
% ways. The explicit \newpage fixes the page break in both files.
\documentclass[10pt]{article}
\usepackage{precalculus-article}
\usepackage{precalculus-boxes}
\newcommand{\TallMath}[1]{$\displaystyle #1\rule[-1.4em]{0pt}{3.2em}$}

\begin{document}

\pageheader{Unit 2, Lesson 2.1}{AP Practice}

\vspace{0.12in}

\boxguard[8]
\begin{remindbox}
\small
\textbf{Extra credit --- optional.} These two pages are written in the style of the AP
Precalculus exam. Your graded homework is the last section of this packet; do it first. Every
answer choice below is a mistake someone really makes, so read all four before you choose.
\end{remindbox}

\vspace{0.12in}

\begin{headlinebox}{goldbg}
\small
\textbf{Section I --- Multiple Choice.} Circle the letter of the single best answer. No calculator.
\end{headlinebox}

\vspace{0.06in}

\begin{enumerate}[leftmargin=1.9em, itemsep=14pt]

\item Which of the following is equivalent to $(2 - i)(3 + 4i)$?

      \smallskip
      \textbf{(A)} $2 + 5i$ \hspace{2.2em} \textbf{(B)} $6 + 5i$ \hspace{2.2em} \textbf{(C)} $10 + 5i$ \hspace{2.2em} \textbf{(D)} $6 - 4i$

\item What is the value of $\sqrt{-4}\cdot\sqrt{-25}$?

      \smallskip
      \textbf{(A)} $10$ \hspace{2.2em} \textbf{(B)} $-10$ \hspace{2.2em} \textbf{(C)} $10i$ \hspace{2.2em} \textbf{(D)} $-10i$

\item Which of the following gives all solutions of $x^2 + 81 = 0$?
      \begin{enumerate}[leftmargin=2.6em, itemsep=2pt]
        \item[(A)] $x = 9$ and $x = -9$
        \item[(B)] $x = 9i$ only
        \item[(C)] There are no solutions of any kind.
        \item[(D)] $x = 9i$ and $x = -9i$
      \end{enumerate}

\item \begin{minipage}[t]{0.56\linewidth}
      The graph of a parent function is shown. The marked points are $(-1,-1)$, $(0,0)$,
      $(1,1)$, and $(8,2)$. Which of the following is its equation?
      \begin{enumerate}[leftmargin=2.6em, itemsep=3pt]
        \item[(A)] $y = \sqrt[3]{x}$
        \item[(B)] $y = x^3$
        \item[(C)] $y = \sqrt{x}$
        \item[(D)] $y = \dfrac{1}{x}$
      \end{enumerate}
      \end{minipage}\hfill
      \begin{minipage}[t]{0.38\linewidth}
      \vspace{-4pt}
      \centering
      \begin{tikzpicture}[x=0.3cm, y=0.55cm, baseline=(current bounding box.north)]
        \draw[linegray, very thin, step=1] (-9,-3) grid (9,3);
        \draw[->, thick] (-9,0) -- (9.6,0) node[right] {\scriptsize $x$};
        \draw[->, thick] (0,-3) -- (0,3.5) node[above] {\scriptsize $y$};
        \foreach \x in {-8,-4,4,8} \node[below, font=\tiny] at (\x,0) {\x};
        \foreach \y in {-2,-1,1,2} \node[left, font=\tiny] at (0,\y) {\y};
        \draw[very thick, plum] plot[domain=-2.08:2.08, samples=60, smooth] ({\x*\x*\x},\x);
        \foreach \p in {(-1,-1),(0,0),(1,1),(8,2)} \fill[plum] \p circle (2pt);
      \end{tikzpicture}
      \end{minipage}

\item Which piecewise function is equal to $|x - 5|$ for every real number $x$?

      \smallskip
      \renewcommand{\arraystretch}{1.3}
      \begin{tabularx}{\linewidth}{@{}XX@{}}
      \textbf{(A)} $\begin{cases} x - 5 & x < 5 \\ -(x - 5) & x \ge 5 \end{cases}$ &
      \textbf{(B)} $\begin{cases} 5 - x & x < 5 \\ x - 5 & x \ge 5 \end{cases}$ \\[14pt]
      \textbf{(C)} $\begin{cases} -x - 5 & x < 5 \\ x - 5 & x \ge 5 \end{cases}$ &
      \textbf{(D)} $\begin{cases} -(x - 5) & x < -5 \\ x - 5 & x \ge -5 \end{cases}$ \\
      \end{tabularx}

\end{enumerate}

\newpage

\begin{headlinebox}{goldbg}
\small
\textbf{Section II --- Free Response.} Show your reasoning. Answers about the part must be
written \emph{in context}, with units --- that is what earns credit on the AP exam.
\end{headlinebox}

\vspace{0.08in}

\begin{enumerate}[leftmargin=1.9em, itemsep=10pt, start=6]

\item \begin{minipage}[t]{0.53\linewidth}
      A machine cuts metal rods that are supposed to be 40 millimeters long. For a rod that is
      actually $L$ millimeters long, the size of the cutting error, in millimeters, is
      \[ E(L) = |L - 40|. \]
      The graph of $E$ for rods from 36 to 44 millimeters is shown.
      \end{minipage}\hfill
      \begin{minipage}[t]{0.40\linewidth}
      \vspace{-2pt}
      \centering
      \begin{tikzpicture}[x=0.52cm, y=0.6cm, baseline=(current bounding box.north)]
        \draw[linegray, very thin, step=1] (36,0) grid (44,4);
        \draw[->, thick] (36,0) -- (44.6,0) node[right] {\scriptsize $L$ (mm)};
        \draw[->, thick] (36,0) -- (36,4.6) node[above] {\scriptsize $E(L)$ (mm)};
        \foreach \x in {36,38,40,42,44} \node[below, font=\tiny] at (\x,0) {\x};
        \foreach \y in {1,2,3,4} \node[left, font=\tiny] at (36,\y) {\y};
        \draw[very thick, plum] (36,4) -- (40,0) -- (44,4);
        \fill[plum] (40,0) circle (2.2pt);
      \end{tikzpicture}
      \end{minipage}

      \textbf{(a)} Find $E(37)$ and $E(42)$. Interpret $E(37)$ in the context of the problem,
      with units.
      \writespace{2.4cm}{}

      \textbf{(b)} Write $E$ as a piecewise function with two rules. State where the break
      between the rules is, and why it is there.
      \writespace{2.4cm}{}

      \textbf{(c)} A technician says the rule for rods shorter than 40 mm must be wrong,
      ``because it has a minus sign in front, so it gives a negative error.'' Use $L = 37$ to
      explain why the technician is wrong.
      \writespace{2.4cm}{}

      \textbf{(d)} The graph of $E$ is the graph of a parent function, moved. Name the parent
      function and state its range. Then explain, in context, why $E(L)$ can never be negative.
      \writespace{2.2cm}{}

      \textbf{(e)} A rod is accepted when its error is at most 1.5 mm. Use each rule from
      part (b) to find the shortest and the longest rod that is accepted.
      \writespace{2.8cm}{}

\end{enumerate}

\end{document}
